
Les relations à prouver se vérifiant ponctuellement sur $\Spec(A)$, on est ramené au cas $A$ local, avec $u(E)\subset\rad(A)$. Les $J^n/J^{n+1}$ sont alors libres sur $B$, et les relations $\Tor_i^A(J^n/J^{n+1},A')=0$ pour $i\geq1$ équivalent aux relations $\Tor_i^A(B,A')=0$ pour $i\geq1$. Or on peut calculer les $\Tor_i^A$ en utilisant la résolution projective $K_\bullet(u)$ de $B$; ce sont les groupes d'homologie du complexe $K_\bullet(u)\otimes_AA'=K_\bullet(u')$. Comme on suppose $u'$ régulier, ce dernier est acyclique en degrés $>0$, d'où les relations $\Tor_i^A(J^n/J^{n+1},A')=0$, pour $i\geq1$.

On en déduit les relations $\Tor_i^A(A/J^n,A')=0$ pour $i\geq1$, en procédant par récurrence sur $n$ grâce à la suite exacte des Tor associée à la suite exacte:
\[
 0\longrightarrow J^n/J^{n+1}\longrightarrow A/J^{n+1}\longrightarrow A/J^n\longrightarrow0.
\]
Comme pour tout $k$ on a $J^kA'=J'^k$, on tire de cette suite exacte et de $\Tor_1^A(A/J^n,A')=0$ les relations $J^n/J^{n+1}\otimes_AA'\xrightarrow{\sim}J'^n/J'^{n+1}$. De la suite exacte:
\[
 0\longrightarrow J^n\longrightarrow A\longrightarrow A/J^n\longrightarrow0,
\]
et de $\Tor_1^A(A/J^n,A')=0$, on déduit les relations $J^n\otimes_AA'\simeq J'^n$. Enfin, de la suite exacte des Tor associée à la suite exacte:
\[
 0\longrightarrow J^{n+1}\longrightarrow J^n\longrightarrow J^n/J^{n+1}\longrightarrow0,
\]
on déduit les relations $\Tor_i^A(J^n,A')=0$ pour $i\geq1$, par récurrence sur $n$, le cas $n=0$ étant évident.

b) On va montrer maintenant que les $J^n/J^{n+1}$ sont projectifs sur $B=A/J$. La question étant locale sur $\Spec(A)$, on peut supposer $E$ libre, soit $E=(A)^d$, de sorte que $u$ est défini par une suite $f_1,\ldots,f_d$ d'éléments de $A$. Posons $A_0=\mathbb Z[T_1,\ldots,T_d]$ et soit $A_0\longrightarrow A$ l'homomorphisme qui applique
